% OpenStax Calculus Volume 3 -- Section 2.3 Exercises: The Dot Product
% Exercises reproduced from OpenStax, licensed under CC BY 4.0.
% Source: https://openstax.org/books/calculus-volume-3/pages/2-3-the-dot-product
\documentclass{exercises}
\title{OpenStax Calculus Volume 3}
\subtitle{Section 2.3 Exercises: The Dot Product}
\sourceurl{https://openstax.org/books/calculus-volume-3/pages/2-3-the-dot-product}
\author{}
\date{}

\begin{document}
\maketitle


For the following exercises, the vectors $\mathbf{u}$ and $\mathbf{v}$ are given. Calculate the dot product $\mathbf{u}\cdot \mathbf{v}$.

\begin{enumerate}[start=1]
\item $\mathbf{u}=\langle3,0\rangle$, $\mathbf{v}=\langle2,2\rangle$
\item $\mathbf{u}=\langle3,-4\rangle$, $\mathbf{v}=\langle4,3\rangle$
\item $\mathbf{u}=\langle2,2,-1\rangle$, $\mathbf{v}=\langle-1,2,2\rangle$
\item $\mathbf{u}=\langle4,5,-6\rangle$, $\mathbf{v}=\langle0,-2,-3\rangle$
\end{enumerate}

For the following exercises, the vectors $\mathbf{a}$, $\mathbf{b}$, and $\mathbf{c}$ are given. Determine the vectors $(\mathbf{a}\cdot \mathbf{b})\mathbf{c}$ and $(\mathbf{a}\cdot \mathbf{c})\mathbf{b}$. Express the vectors in component form.

\begin{enumerate}[resume]
\item $\mathbf{a}=\langle2,0,-3\rangle$, $\mathbf{b}=\langle-4,-7,1\rangle$, $\mathbf{c}=\langle1,1,-1\rangle$
\item $\mathbf{a}=\langle0,1,2\rangle$, $\mathbf{b}=\langle-1,0,1\rangle$, $\mathbf{c}=\langle1,0,-1\rangle$
\item $\mathbf{a}=\mathbf{i}+\mathbf{j}$, $\mathbf{b}=\mathbf{i}-\mathbf{k}$, $\mathbf{c}=\mathbf{i}-2\mathbf{k}$
\item $\mathbf{a}=\mathbf{i}-\mathbf{j}+\mathbf{k}$, $\mathbf{b}=\mathbf{j}+3\mathbf{k}$, $\mathbf{c}=-\mathbf{i}+2\mathbf{j}-4\mathbf{k}$
\end{enumerate}

For the following exercises, the two-dimensional vectors $\mathbf{a}$ and $\mathbf{b}$ are given.


(a) Find the measure of the angle $\theta$ between $\mathbf{a}$ and $\mathbf{b}$. Express the answer in radians rounded to two decimal places, if it is not possible to express it exactly.


(b) Is $\theta$ an acute angle?

\begin{enumerate}[resume]
\item \techrequired $\mathbf{a}=\langle3,-1\rangle$, $\mathbf{b}=\langle-4,0\rangle$
\item \techrequired $\mathbf{a}=\langle2,1\rangle$, $\mathbf{b}=\langle-1,3\rangle$
\item $\mathbf{u}=3\mathbf{i}$, $\mathbf{v}=4\mathbf{i}+4\mathbf{j}$
\item $\mathbf{u}=5\mathbf{i}$, $\mathbf{v}=-6\mathbf{i}+6\mathbf{j}$
\end{enumerate}

For the following exercises, find the measure of the angle between the three-dimensional vectors $\mathbf{a}$ and $\mathbf{b}$. Express the answer in radians rounded to two decimal places, if it is not possible to express it exactly.

\begin{enumerate}[resume]
\item $\mathbf{a}=\langle3,-1,2\rangle$, $\mathbf{b}=\langle1,-1,-2\rangle$
\item $\mathbf{a}=\langle0,-1,-3\rangle$, $\mathbf{b}=\langle2,3,-1\rangle$
\item $\mathbf{a}=\mathbf{i}+\mathbf{j}$, $\mathbf{b}=\mathbf{j}-\mathbf{k}$
\item $\mathbf{a}=\mathbf{i}-2\mathbf{j}+\mathbf{k}$, $\mathbf{b}=\mathbf{i}+\mathbf{j}-2\mathbf{k}$
\item \techrequired $\mathbf{a}=3\mathbf{i}-\mathbf{j}-2\mathbf{k}$, $\mathbf{b}=\mathbf{v}+\mathbf{w}$, where $\mathbf{v}=-2\mathbf{i}-3\mathbf{j}+2\mathbf{k}$ and $\mathbf{w}=\mathbf{i}+2\mathbf{k}$
\item \techrequired $\mathbf{a}=3\mathbf{i}-\mathbf{j}+2\mathbf{k}$, $\mathbf{b}=\mathbf{v}-\mathbf{w}$, where $\mathbf{v}=2\mathbf{i}+\mathbf{j}+4\mathbf{k}$ and $\mathbf{w}=6\mathbf{i}+\mathbf{j}+2\mathbf{k}$
\end{enumerate}

For the following exercises determine whether the given vectors are orthogonal.

\begin{enumerate}[resume]
\item $\mathbf{a}=\langle x,y\rangle$, $\mathbf{b}=\langle-y,x\rangle$, where $x$ and $y$ are nonzero real numbers
\item $\mathbf{a}=\langle x,x\rangle$, $\mathbf{b}=\langle-y,y\rangle$, where $x$ and $y$ are nonzero real numbers
\item $\mathbf{a}=3\mathbf{i}-\mathbf{j}-2\mathbf{k}$, $\mathbf{b}=-2\mathbf{i}-3\mathbf{j}+\mathbf{k}$
\item $\mathbf{a}=\mathbf{i}-\mathbf{j}$, $\mathbf{b}=7\mathbf{i}+2\mathbf{j}-\mathbf{k}$
\item Find all two-dimensional vectors $\mathbf{a}$ orthogonal to vector $\mathbf{b}=\langle3,4\rangle$. Express the answer in component form.
\item Find all two-dimensional vectors $\mathbf{a}$ orthogonal to vector $\mathbf{b}=\langle5,-6\rangle$. Express the answer by using standard unit vectors.
\item Determine all three-dimensional vectors $\mathbf{u}$ orthogonal to vector $\mathbf{v}=\langle1,1,0\rangle$. Express the answer by using standard unit vectors.
\item Determine all three-dimensional vectors $\mathbf{u}$ orthogonal to vector $\mathbf{v}=\mathbf{i}-\mathbf{j}-\mathbf{k}$. Express the answer in component form.
\item Determine the real number $\alpha$ such that vectors $\mathbf{a}=2\mathbf{i}+3\mathbf{j}$ and $\mathbf{b}=9\mathbf{i}+\alpha \mathbf{j}$ are orthogonal.
\item Determine the real number $\alpha$ such that vectors $\mathbf{a}=-3\mathbf{i}+2\mathbf{j}$ and $\mathbf{b}=2\mathbf{i}+\alpha \mathbf{j}$ are orthogonal.
\item \techrequired Consider the points $P(4,5)$ and $Q(5,-7)$.
\begin{enumerate}[label=(\alph*)]
  \item Determine vectors $\overset{\to}{OP}$ and $\overset{\to}{OQ}$. Express the answer by using standard unit vectors.
  \item Determine the measure of angle O in triangle OPQ. Express the answer in degrees rounded to two decimal places.
\end{enumerate}
\item \techrequired Consider points $A(1,1)$, $B(2,-7)$, and $C(6,3)$.
\begin{enumerate}[label=(\alph*)]
  \item Determine vectors $\overset{\to}{BA}$ and $\overset{\to}{BC}$. Express the answer in component form.
  \item Determine the measure of angle B in triangle ABC. Express the answer in degrees rounded to two decimal places.
\end{enumerate}
\item Determine the measure of angle A in triangle ABC, where $A(1,1,8)$, $B(4,-3,-4)$, and $C(-3,1,5)$. Express your answer in degrees rounded to two decimal places.
\item Consider points $P(3,7,-2)$ and $Q(1,1,-3)$. Determine the angle between vectors $\overset{\to}{OP}$ and $\overset{\to}{OQ}$. Express the answer in degrees rounded to two decimal places.
\end{enumerate}

For the following exercises, determine which (if any) pairs of the following vectors are orthogonal.

\begin{enumerate}[resume]
\item $\mathbf{u}=\langle3,7,-2\rangle$, $\mathbf{v}=\langle5,-3,-3\rangle$, $\mathbf{w}=\langle0,1,-1\rangle$
\item $\mathbf{u}=\mathbf{i}-\mathbf{k}$, $\mathbf{v}=5\mathbf{j}-5\mathbf{k}$, $\mathbf{w}=10\mathbf{j}$
\item Use vectors to show that a parallelogram with equal diagonals is a rectangle.
\item Use vectors to show that the diagonals of a rhombus are perpendicular.
\item Show that $\mathbf{u}\cdot(\mathbf{v}+\mathbf{w})=\mathbf{u}\cdot \mathbf{v}+\mathbf{u}\cdot \mathbf{w}$ is true for any vectors $\mathbf{u}$, $\mathbf{v}$, and $\mathbf{w}$.
\item Verify the identity $\mathbf{u}\cdot(\mathbf{v}+\mathbf{w})=\mathbf{u}\cdot \mathbf{v}+\mathbf{u}\cdot \mathbf{w}$ for vectors $\mathbf{u}=\langle1,0,4\rangle$, $\mathbf{v}=\langle-2,3,5\rangle$, and $\mathbf{w}=\langle4,-2,6\rangle$.
\end{enumerate}

For the following problems, the vector $\mathbf{u}$ is given.


(a) Find the direction cosines for the vector $\mathbf{u}$.


(b) Find the direction angles for the vector $\mathbf{u}$ expressed in degrees. (Round the answer to the nearest integer.)

\begin{enumerate}[resume]
\item $\mathbf{u}=\langle2,2,1\rangle$
\item $\mathbf{u}=\mathbf{i}-2\mathbf{j}+2\mathbf{k}$
\item $\mathbf{u}=\langle-1,5,2\rangle$
\item $\mathbf{u}=\langle2,3,4\rangle$
\item Consider $\mathbf{u}=\langle a,b,c\rangle$ a nonzero three-dimensional vector. Let $\cos \alpha$, $\cos \beta$, and $\cos \gamma$ be the direction cosines of $\mathbf{u}$. Show that $\cos^{2}\alpha+\cos^{2}\beta+\cos^{2}\gamma=1$.
\item Determine the direction cosines of vector $\mathbf{u}=\mathbf{i}+2\mathbf{j}+2\mathbf{k}$ and show they satisfy $\cos^{2}\alpha+\cos^{2}\beta+\cos^{2}\gamma=1$.
\end{enumerate}

For the following exercises, the vectors $\mathbf{u}$ and $\mathbf{v}$ are given.


(a) Find the vector projection $\mathbf{w}=\operatorname{proj}_{\mathbf{u}}\mathbf{v}$ of vector $\mathbf{v}$ onto vector $\mathbf{u}$. Express your answer in component form.


(b) Find the scalar projection $\operatorname{comp}_{\mathbf{u}}\mathbf{v}$ of vector $\mathbf{v}$ onto vector $\mathbf{u}$.

\begin{enumerate}[resume]
\item $\mathbf{u}=5\mathbf{i}+2\mathbf{j}$, $\mathbf{v}=2\mathbf{i}+3\mathbf{j}$
\item $\mathbf{u}=\langle-4,7\rangle$, $\mathbf{v}=\langle3,5\rangle$
\item $\mathbf{u}=3\mathbf{i}+2\mathbf{k}$, $\mathbf{v}=2\mathbf{j}+4\mathbf{k}$
\item $\mathbf{u}=\langle4,4,0\rangle$, $\mathbf{v}=\langle0,4,1\rangle$
\item Consider the vectors $\mathbf{u}=4\mathbf{i}-3\mathbf{j}$ and $\mathbf{v}=3\mathbf{i}+2\mathbf{j}$.
\begin{enumerate}[label=(\alph*)]
  \item Find the component form of vector $\mathbf{w}=\operatorname{proj}_{\mathbf{u}}\mathbf{v}$ that represents the projection of $\mathbf{v}$ onto $\mathbf{u}$.
  \item Write the decomposition $\mathbf{v}=\mathbf{w}+\mathbf{q}$ of vector $\mathbf{v}$ into the orthogonal components $\mathbf{w}$ and $\mathbf{q}$, where $\mathbf{w}$ is the projection of $\mathbf{v}$ onto $\mathbf{u}$ and $\mathbf{q}$ is a vector orthogonal to the direction of $\mathbf{u}$.
\end{enumerate}
\item Consider vectors $\mathbf{u}=2\mathbf{i}+4\mathbf{j}$ and $\mathbf{v}=4\mathbf{j}+2\mathbf{k}$.
\begin{enumerate}[label=(\alph*)]
  \item Find the component form of vector $\mathbf{w}=\operatorname{proj}_{\mathbf{u}}\mathbf{v}$ that represents the projection of $\mathbf{v}$ onto $\mathbf{u}$.
  \item Write the decomposition $\mathbf{v}=\mathbf{w}+\mathbf{q}$ of vector $\mathbf{v}$ into the orthogonal components $\mathbf{w}$ and $\mathbf{q}$, where $\mathbf{w}$ is the projection of $\mathbf{v}$ onto $\mathbf{u}$ and $\mathbf{q}$ is a vector orthogonal to the direction of $\mathbf{u}$.
\end{enumerate}
\item A methane molecule has a carbon atom situated at the origin and four hydrogen atoms located at points $P(1,1,-1),Q(1,-1,1),R(-1,1,1)$, and $S(-1,-1,-1)$ (see figure). \\ \textit{[Figure: This figure is the 3-dimensional coordinate system. There are four points plotted. The first point is labeled ``P(1, 1, -1),'' the second point is labeled ``Q(1, -1, 1),'' the third point is labeled ``R(-1, 1, 1),'' and the fourth point is labeled ``S(-1, -1, -1).'' There are line segments from Q to P, P to R and R to Q. There are also two vectors in standard position. The first has terminal point of R and the second has terminal point of S. The angle between them is represented with an arc.]}
\begin{enumerate}[label=(\alph*)]
  \item Find the distance between the hydrogen atoms located at P and R.
  \item Find the angle between vectors $\overset{\to}{OS}$ and $\overset{\to}{OR}$ that connect the carbon atom with the hydrogen atoms located at S and R, which is also called the bond angle. Express the answer in degrees rounded to two decimal places.
\end{enumerate}
\item \techrequired Find the vectors that join the center of a clock to the hours 1:00, 2:00, and 3:00. Assume the clock is circular with a radius of 1 unit.
\item Find the work done by force $\mathbf{F}=\langle5,6,-2\rangle$ (measured in Newtons) that moves a particle from point $P(3,-1,0)$ to point $Q(2,3,1)$ along a straight line (the distance is measured in meters).
\item \techrequired A sled is pulled by exerting a force of 100 N on a rope that makes an angle of $25^\circ$ with the horizontal. Find the work done in pulling the sled 40 m. (Round the answer to one decimal place.)
\item \techrequired A father is pulling his son on a sled at an angle of $20^\circ$ with the horizontal with a force of 25 lb (see the following image). He pulls the sled in a straight path of 50 ft. How much work was done by the man pulling the sled? (Round the answer to the nearest integer.) \\ \textit{[Figure: This figure is an image of a person pulling a child on a sled. The rope for pulling the sled is represented by a vector and labeled ``25 lb.'' There is an angle between the rope vector and the horizontal ground of 20 degrees.]}
\item \techrequired A car is towed using a force of 1600 N. The rope used to pull the car makes an angle of $25^\circ$ with the horizontal. Find the work done in towing the car 2 km. Express the answer in joules $(1\text{J}=1\text{N}\cdot \text{m})$ rounded to the nearest integer.
\item \techrequired A boat sails north aided by a wind blowing in a direction of $\text{N}30^\circ\text{E}$ with a magnitude of 500 lb. How much work is performed by the wind as the boat moves 100 ft? (Round the answer to two decimal places.)
\item Vector $\mathbf{p}=\langle150,225,375\rangle$ represents the price of certain models of bicycles sold by a bicycle shop. Vector $\mathbf{n}=\langle10,7,9\rangle$ represents the number of bicycles sold of each model, respectively. Compute the dot product $\mathbf{p}\cdot \mathbf{n}$ and state its meaning.
\item \techrequired Two forces $\mathbf{F}_{1}$ and $\mathbf{F}_{2}$ are represented by vectors with initial points that are at the origin. The first force has a magnitude of 20 lb and the terminal point of the vector is point $P(1,1,0)$. The second force has a magnitude of 40 lb and the terminal point of its vector is point $Q(0,1,1)$. Let $\mathbf{F}$ be the resultant force of forces $\mathbf{F}_{1}$ and $\mathbf{F}_{2}$.
\begin{enumerate}[label=(\alph*)]
  \item Find the magnitude of $\mathbf{F}$. (Round the answer to one decimal place.)
  \item Find the direction angles of $\mathbf{F}$. (Express the answer in degrees rounded to one decimal place.)
\end{enumerate}
\item \techrequired Consider $\mathbf{r}(t)=\langle \cos t,\sin t,2t\rangle$ the position vector of a particle at time $t\in[0,30]$, where the components of $\mathbf{r}$ are expressed in centimeters and time in seconds. Let $\overset{\to}{OP}$ be the position vector of the particle after 1 sec.
\begin{enumerate}[label=(\alph*)]
  \item Show that all vectors $\overset{\to}{PQ}$, where $Q(x,y,z)$ is an arbitrary point, orthogonal to the instantaneous velocity vector $\mathbf{v}(1)$ of the particle after 1 sec, can be expressed as $\overset{\to}{PQ}=\langle x-\cos 1,y-\sin 1,z-2\rangle$, where $x\sin 1-y\cos 1-2z+4=0$. The set of point Q describes a plane called the normal plane to the path of the particle at point P.
  \item Use a CAS to visualize the instantaneous velocity vector and the normal plane at point P along with the path of the particle.
\end{enumerate}
\end{enumerate}

\end{document}
